\section{Codificación y formatos de vídeo}

La codificación de vídeo es el proceso de reducir el volumen de datos que representa
una secuencia de imágenes en movimiento mediante la eliminación de redundancia espacial
(dentro de cada fotograma) y temporal (entre fotogramas consecutivos). La elección del
\textit{codec} y del contenedor condiciona directamente la calidad de imagen, la
latencia de procesamiento, el ancho de banda necesario y la compatibilidad con
plataformas de distribución.

\subsection{Codecs relevantes}

\begin{itemize}
  \item \textbf{H.264 / AVC} (\textit{Advanced Video Coding}, ITU-T H.264 / ISO 14496-10):
  el codec más extendido en la actualidad. Ofrece buena calidad a tasas de bits moderadas
  y es soportado por prácticamente toda la infraestructura de streaming existente
  (YouTube, Twitch, hardware encoders). Es el codec empleado habitualmente en la
  salida de streaming de Voctomix mediante FFmpeg.
  \item \textbf{H.265 / HEVC} (\textit{High Efficiency Video Coding}): sucesor de H.264
  que reduce el bitrate necesario aproximadamente a la mitad para la misma calidad
  perceptual. Su adopción está limitada por el coste computacional de la codificación
  y por restricciones de licencias en algunas plataformas.
  \item \textbf{VP8 / VP9}: codecs libres de royalties desarrollados por Google, empleados
  principalmente en WebRTC (VP8) y YouTube (VP9). Ofrecen eficiencia comparable a H.264/H.265
  sin restricciones de licencia.
  \item \textbf{AV1}: codec abierto desarrollado por la Alliance for Open Media. Es el
  sucesor de VP9 y el más eficiente de los mencionados, aunque su adopción en tiempo
  real aún está limitada por la demanda de cómputo del encoder.
\end{itemize}

\subsection{Formatos de vídeo sin comprimir}

Internamente, Voctomix opera sobre vídeo \textbf{sin comprimir} para evitar artefactos
acumulativos en el proceso de composición. Los formatos más relevantes son:

\begin{itemize}
  \item \textbf{UYVY} (también denominado YUY2 en variante packed): formato de crominancia
  submuestreada 4:2:2, con 2~bytes por píxel. Es el formato nativo de entrada y salida
  de voctocore.
  \item \textbf{I420} (YUV 4:2:0 planar): formato habitual en la entrada de encoders
  de H.264, con menor resolución cromática que UYVY pero suficiente para la mayoría
  de contenidos.
\end{itemize}

\subsection{Contenedores}

El contenedor define cómo se encapsulan y multiplexan los flujos de vídeo y audio:

\begin{itemize}
  \item \textbf{Matroska (MKV)}: contenedor flexible y extensible, empleado por
  Voctomix para el transporte interno de señal sin comprimir sobre TCP.
  \item \textbf{MPEG-TS} (\textit{Transport Stream}): estándar de broadcast diseñado
  para transmisión con pérdidas. Es el contenedor habitual en emisión por satélite,
  TDT y algunos protocolos de streaming como HLS.
  \item \textbf{MP4 / ISOBMFF}: contenedor de referencia para distribución de ficheros
  y streaming HTTP (DASH, HLS fragmentado).
  \item \textbf{FLV}: contenedor legado asociado a RTMP, aún utilizado en la mayoría
  de plataformas de streaming en directo.
\end{itemize}

\subsection{FFmpeg como herramienta de referencia}

FFmpeg es la herramienta de línea de comandos de código abierto más utilizada para el
procesamiento de audio y vídeo en entornos Linux. Incorpora prácticamente todos los
codecs, contenedores y filtros descritos en este capítulo, actuando como capa de
abstracción sobre las bibliotecas subyacentes (libx264, libx265, libvpx, libaom).

En el contexto de Voctomix 2.0, FFmpeg desempeña un doble papel: como
\textbf{generador de fuentes de prueba} mediante los scripts incluidos en
\texttt{example-scripts/ffmpeg/} (barras de color, \textit{test card} con reloj),
y como \textbf{encoder de salida} que convierte la mezcla en bruto del puerto 11000
en un flujo H.264/AAC apto para plataformas como YouTube Live o Twitch.
